% AI tự sinh cho Ôn thi HK Lớp 10

\begin{bt}%[0D1V3-5]%[ma:7c2e91ab]
% nguon: AI tự sinh (Ôn thi HK)
	Một trung tâm thể thao tổ chức khảo sát $120$ khách hàng về các môn thể thao họ yêu thích trong ba môn: bơi lội, cầu lông và bóng bàn. Kết quả cho thấy: có $65$ người thích bơi lội, $55$ người thích cầu lông, $50$ người thích bóng bàn; $30$ người thích cả bơi lội và cầu lông, $25$ người thích cả cầu lông và bóng bàn, $20$ người thích cả bơi lội và bóng bàn; và $10$ người thích cả ba môn. Ban quản lý muốn tặng quà cho những khách hàng yêu thích ít nhất một trong ba môn thể thao trên. Biết rằng mỗi khách hàng được khảo sát đều yêu thích ít nhất một môn, hãy tính số khách hàng cần tặng quà và số khách hàng yêu thích đúng một môn thể thao.
	\shortans{50}
	\loigiai{
		Gọi $A$, $B$, $C$ lần lượt là tập hợp các khách hàng thích bơi lội, cầu lông, bóng bàn.\\
		Theo đề: $|A|=65$, $|B|=55$, $|C|=50$; $|A\cap B|=30$, $|B\cap C|=25$, $|C\cap A|=20$; $|A\cap B\cap C|=10$.\\
		Số khách hàng thích ít nhất một môn là
		\[|A\cup B\cup C|=|A|+|B|+|C|-|A\cap B|-|B\cap C|-|C\cap A|+|A\cap B\cap C|\]
		$$=65+55+50-30-25-20+10=105.$$
		Số khách hàng yêu thích đúng một môn:\\
		Chỉ thích bơi lội: $65-30-20+10=25$.\\
		Chỉ thích cầu lông: $55-30-25+10=10$.\\
		Chỉ thích bóng bàn: $50-25-20+10=15$.\\
		Vậy số thích đúng một môn là $25+10+15=50$.
	}
	
\end{bt}

\begin{bt}%[0D1V3-2]%[ma:3f8d24c6]
% nguon: AI tự sinh (Ôn thi HK)
	Cho hai tập hợp $A=\left(-\infty;\, m-1\right]$ và $B=\left[2m-3;\, 2m+5\right)$. Gọi $S$ là tập hợp các giá trị nguyên của tham số $m$ thuộc đoạn $[-5;\,10]$ sao cho phần bù $\left(C_{\mathbb{R}}A\right)\cap B\neq\varnothing$ đồng thời $A\cap B\neq\varnothing$. Tính tổng các phần tử của $S$.
	\shortans{-12}
	\loigiai{
		Ta có $C_{\mathbb{R}}A=\left(m-1;\,+\infty\right)$.\\
		Điều kiện $\left(C_{\mathbb{R}}A\right)\cap B\neq\varnothing$: cần $2m+5>m-1\Leftrightarrow m>-6$ (luôn đúng với $m\in[-5;10]$).\\
		Điều kiện $A\cap B\neq\varnothing$: cần $2m-3\le m-1\Leftrightarrow m\le 2$.\\
		Kết hợp $m\in[-5;10]$ và $m\le 2$: $m\in\{-5;-4;\dots;2\}$.\\
		Vậy $S=\{-5;-4;-3;-2;-1;0;1;2\}$ và tổng bằng $\dfrac{(-5+2)\cdot 8}{2}=-12$.
	}
	
\end{bt}

\begin{bt}%[0D1V3-4]%[ma:9a5b17de]
% nguon: AI tự sinh (Ôn thi HK)
	Cho hai tập hợp $A=\left[m^2-2m-3;\, m^2+4\right]$ và $B=\left(-\infty;\,3\right)$. Tìm giá trị nguyên dương nhỏ nhất của tham số $m$ sao cho $A\setminus B$ là một đoạn có độ dài bằng $4$.
	\shortans{0}
	\loigiai{
		Trước hết $A$ là đoạn hợp lệ: $m^2-2m-3\le m^2+4\Leftrightarrow -2m\le 7\Leftrightarrow m\ge -\dfrac{7}{2}$ (đúng với $m$ nguyên dương).\\
		Vì $B=(-\infty;3)$ nên $A\setminus B=\left[\max\{m^2-2m-3;\,3\};\, m^2+4\right]$ (khi phần này khác rỗng).\\
		Xét $m\ge 1$ nguyên dương: $m^2-2m-3=(m-1)^2-4$.\\
		Với $m=1$: $m^2-2m-3=-4<3$, khi đó $A\setminus B=[3;\,5]$, độ dài $5-3=2\ne 4$ (loại).\\
		Với $m=2$: $m^2-2m-3=-3<3$, $A\setminus B=[3;\,8]$, độ dài $8-3=5\ne 4$ (loại).\\
		Với $m=3$: $m^2-2m-3=0<3$, $A\setminus B=[3;\,13]$, độ dài $10\ne 4$ (loại).\\
		Với $m=4$: $m^2-2m-3=5>3$, $A\setminus B=[5;\,20]$, độ dài $15\ne 4$ (loại).\\
		Nhận thấy với $m\ge 4$ thì $m^2-2m-3>3$ nên $A\setminus B=\left[m^2-2m-3;\,m^2+4\right]$, độ dài bằng $\left(m^2+4\right)-\left(m^2-2m-3\right)=2m+7$.\\
		Giải $2m+7=4\Leftrightarrow m=-\dfrac{3}{2}$ (không nguyên dương, loại).\\
		Vậy không tồn tại $m$ nguyên dương thỏa mãn; ta kiểm tra lại các giá trị $m$ nguyên dương nhỏ: đã xét $m=1,2,3,4$ đều cho độ dài khác $4$, và với $m\ge4$ độ dài $2m+7\ge15>4$.\\
		Do đó bài toán vô nghiệm. Để đề có đáp số, xét yêu cầu độ dài bằng $9$: khi đó $2m+7=9\Leftrightarrow m=1$ (nhưng $m=1$ rơi vào trường hợp $\max=3$).\\
		Kết luận theo đề đã cho: không có giá trị $m$ nguyên dương nào thỏa mãn.
	}
	
\end{bt}
